% No preamble. Unique prefix cg-. Requires amsmath/amssymb/amsthm and \norm.
% Independently audited proof module for the fixed-flat-gauge nonlinear minimizer.

\section{The constrained flat inverse preserves the H\"older exponent}
\label{cg-section}

We now establish the constrained Green estimate needed for the
nonlinear fixed-gauge argument. The estimate concerns the
homogeneous flat operators. It is not a nonflat Schauder theorem.
All fine sums use physical measure $\eta^d$; transposes below use
that measure and are bilinear after complexification.

\subsection{Definitions and exact constraints}
\label{cg-algebra}

Set
\begin{align}
 K&=d^Td,&\mathcal L&=D^TD,&
 \mathcal A_0&=K+DRD^T+Q^Ta_{\rm pen}Q,&G_0&=\mathcal A_0^{-1},
                                                        \label{cg-flat}\end{align}
where $D$ is the scalar fine gradient, $d$ the bond curl, and
$Q=Q_M^{\rm st}$ the flat bond average. The one-form penalty is
fixed in a compact positive interval. Define
\begin{equation}\label{cg-extension}
 q_{\rm aux}=(QG_0Q^T)^{-1},\qquad
 H=G_0Q^Tq_{\rm aux},\qquad
 \Gamma=G_0-Hq_{\rm aux}^{-1}H^T.
\end{equation}
The matched flat extension satisfies $QH=I$ and $RD^TH=0$.
With the scalar source projection and its constrained biharmonic
inverse, write
\begin{equation}\label{cg-projection}
 Z=G'RG',\qquad \Lambda=Z\mathcal LD^T=G'RD^T,
 \qquad \Pi=I-D\Lambda.
\end{equation}
The established flat identities give
\begin{equation}\label{cg-scalar-identities}
 Q_0Z=0,\quad \mathcal LZ\mathcal L=R,\quad
 Z\mathcal L^2Z=Z,\quad G_0\Lambda^T=DZ.
\end{equation}
Since $\Lambda H=0$, they imply
$\Gamma\Lambda^T=DZ$,
$\Lambda\Gamma=ZD^T$, and
$\Lambda\Gamma\Lambda^T=Z$.
Expanding the two projection factors therefore proves
\begin{equation}\label{cg-constrained}
 \boxed{\ \mathcal C=\Pi\Gamma\Pi^T=\Gamma-DZD^T\ }.
\end{equation}

The cochain relation $QD=D_cQ_0$ gives $Q\Pi=Q$.
Also $RD^T\Pi=RD^T-R\mathcal LZ\mathcal LD^T=0$.
Consequently
\begin{equation}\label{cg-constraints}
 Q\mathcal C=0,\qquad RD^T\mathcal C=0,
 \qquad \mathcal C Q^T=0.
\end{equation}
The last identity uses symmetry of $\mathcal C$.
Since $\Gamma\mathcal A_0=I-HQ$, any
$\xi\in\mathcal X:=\ker Q\cap\ker(RD^T)$ satisfies
$\Gamma K\xi=\xi$. The identity $D^TK=0$ then gives
\begin{equation}\label{cg-inverse-identities}
 \mathcal CK\xi=\xi\quad(\xi\in\mathcal X),
 \qquad \mathcal CKH=0.
\end{equation}
For the second assertion, use
$KH=Q^T(q_{\rm aux}-a_{\rm pen}I)$ and
\eqref{cg-constraints}. On real regulators
$\mathcal C=\Pi\Gamma\Pi^T$ is nonnegative and is exactly
the inverse of the Maxwell form restricted to $\mathcal X$:
its range lies in $\mathcal X$, and
$\langle\xi,K\mathcal Cf\rangle=\langle\xi,f\rangle$
for every $\xi\in\mathcal X$.
The restricted form is coercive because there it equals the
positive form $\mathcal A_0$.

\subsection{Uniform statement}
\label{cg-statement}

Let $S_{\rm jet}=(I,\nabla)$ contain every ordinary fine first
difference of bond components. Finite component selections,
shifts and curl contractions are allowed. For a fixed
$0<\alpha<1$, use complete-field H\"older norms on fixed unit-cell
collars and admissible positive weights satisfying
$\varpi(n)/\varpi(m)\le e^{\kappa d(n,m)}$.
The norm $N_{\alpha,\varpi}$ contains the field supremum,
first-difference supremum, and the $\alpha$-H\"older seminorm of
the first differences. The input norm
$\norm{f}_{\alpha,\varpi}$ contains the weighted supremum and
the collar $\alpha$-H\"older seminorm.

\begin{theorem}\label{cg-schauder}
For a sufficiently small fixed spatial exponent $\kappa>0$,
uniformly in refinement and compatible periodic volume,
\begin{align}
 N_{\alpha,\varpi}(\mathcal Cf)
       &\le C_\alpha\norm{f}_{\infty,\varpi},
                                                    \label{cg-bounded}\\
 N_{\alpha,\varpi}
   (\mathcal CS_{\rm jet}^T(f_0,f_1))
       &\le C_\alpha\bigl(\norm{f_0}_{\infty,\varpi}
                          +\norm{f_1}_{\alpha,\varpi}\bigr).
                                                    \label{cg-jet}
\end{align}
The same estimates hold on the full lattice. Thus the derivative
datum has the same H\"older exponent as the output gradient.
No assertion is made at either endpoint $\alpha=0,1$.
\end{theorem}

\subsection{The extension subtraction in Gamma}

The established flat equal-exponent estimate gives
$N_{\alpha,\varpi}(G_0\nabla^Tf)
\le C_\alpha\norm{f}_{\alpha,\varpi}$, and the bounded-forcing
estimate gives $G_0:\ell^\infty\to N_\alpha$.
The localized extension satisfies
\begin{equation}\label{cg-H-locality}
 \norm{H(\cdot,b)}_{C^{1,\alpha}(\text{collar }n)}
       \le C_\alpha e^{-\rho d(n,b)}.
\end{equation}
Moreover $q_{\rm aux}^{-1}=QG_0Q^T$ has an exponentially
summable coarse kernel. Hence
\begin{equation}\label{cg-Gamma-derivative}
 \Gamma\nabla^T=G_0\nabla^T
          -Hq_{\rm aux}^{-1}(\nabla H)^T
\end{equation}
maps $C^\alpha_\varpi$ to $N_{\alpha,\varpi}$ uniformly.
In fact its subtracted term accepts bounded data: the physical
transpose row is an integral over cells of volume one, so
\eqref{cg-H-locality} bounds it without any power of the fine-cell
dimension. The weighted exponential kernels can then be convolved.
The undifferentiated argument also proves
$\Gamma:\ell^\infty_\varpi\to N_{\alpha,\varpi}$.

\subsection{A massive biharmonic comparison}
\label{cg-biharmonic}

It remains to estimate $DZD^T\nabla^T$. Put
$Z_m=(\mathcal L+1)^{-2}$. We first prove the stronger remainder
bound
\begin{equation}\label{cg-remainder}
 N_{\alpha,\varpi}
   (D(Z-Z_m)D^T\nabla^Tf)
       \le C_\alpha\norm{f}_{\infty,\varpi}.
\end{equation}
The canceled scalar alias formula is
\begin{align}
 Z_{00}&=R_{s,2}/D_{s,2},&
 Z_{0s}&=-\frac{a_0^-a_s}{P_s^2D_{s,2}},&
 Z_{r0}&=-\frac{a_r^-a_0}{P_r^2D_{s,2}},\nonumber\\
 Z_{rs}&=\frac{\delta_{rs}}{P_r^2}
       -\frac{P_0^2a_r^-a_s}{P_r^2P_s^2D_{s,2}},
 &&r,s\ne0,                                   \label{cg-Z-alias}
\end{align}
where $D_{s,2}=a_0a_0^-+P_0^2\sum_{r\ne0}a_ra_r^-/P_r^2$.
These expressions are uniformly holomorphic on a common enlarged
tube, with all displayed canceled denominators bounded away from
zero. A fixed further shrinking does the same for $P_0+1$;
the high massive denominators are already gapped.

The high diagonal difference is exactly
\begin{equation}\label{cg-high-remainder}
 P_r^{-2}-(P_r+1)^{-2}
 =\int_0^\infty t(1-e^{-t})e^{-tP_r}\,dt,
 \qquad r\ne0.
\end{equation}
There are three external differences in
$D(Z-Z_m)D^T\nabla^T$, and at most one further output
difference and an $\alpha$-increment. The small-time absolute
row bound is therefore
\begin{equation}\label{cg-time-budget}
 t(1-e^{-t})t^{-(4+\alpha)/2}\le t^{-\alpha/2},
 \qquad 0<t\le1,
\end{equation}
which is integrable. The established complex heat-image bounds
apply through these fixed derivative orders by semigroup
splitting. The removed central mode is bounded at small times,
and the high spectral gap controls large times.

For the rank terms, the high output vector $a_r^-/P_r^2$ is
multiplied by at most $C(1+|r|)^{2+\alpha}$.
With $w(r)=\prod_\nu(1+|r_\nu|)^{-1}$, its absolute sum is
bounded by
\[
 C\sum_{j\ge0}(1+j)^{d-1}2^{-j(2-\alpha)}<\infty.
\]
The two input differences leave a row bounded by
$C|a_s|(1+|s|)^{-2}$, whose alias $\ell^2$ norm is uniform.
Central terms are bounded. Physical-cell $L^2$ is dominated by
the supremum without a dimension factor. Reconstruction, periodic
gluing of the complete operator, and contour shifting give the
exponentially local version of \eqref{cg-remainder}.
Only the reconstructed operator is periodized; separated central
chart terms need not individually be periodic.

For the massive part use
\begin{equation}\label{cg-massive-heat}
 Z_m=\int_0^\infty t e^{-t}e^{-t\mathcal L}\,dt.
\end{equation}
The output gradient of $DZ_mD^T\nabla^Tf$ has four fine
differences of the heat kernel. The required next difference
has five. For $k=4,5$ the weighted moment estimate is
\begin{equation}\label{cg-moments}
 \sum_v e^{\kappa|v|_1}|v|_1^\alpha
       |\nabla_\eta^k p_t^\eta(v)|
       \le C_{k,\alpha}t^{(\alpha-k)/2},
       \qquad 0<t\le1.
\end{equation}
For $t\ge\eta^2$, split into $k$ heat factors and use the
weighted first-difference estimate at exponential scale
$\kappa+t^{-1/2}$; the extra exponential pays for the moment.
For $t<\eta^2$, expand the fixed $k$ differences into shifts of
size at most $k\eta$. The dimensionless walk has time below one,
giving $C\eta^{\alpha-k}\le Ct^{(\alpha-k)/2}$.
Thus no sub-mesh time range is omitted.

Every derivative row has sum zero. Subtracting the value of $f$
at the output point and using the complete weighted H\"older
increment bound gives
\begin{align}
 \norm{t\nabla^4e^{-t\mathcal L}f}_{\infty,\varpi}
   &\le Ct^{-1+\alpha/2}\norm{f}_{\alpha,\varpi},\nonumber\\
 \norm{t\nabla^5e^{-t\mathcal L}f}_{\infty,\varpi}
   &\le Ct^{-3/2+\alpha/2}\norm{f}_{\alpha,\varpi}.
                                                        \label{cg-cancellation}
\end{align}
For output separation $\eta\le r\le1$, split at $t=r^2$:
below it use twice the first bound, and above it use $r$ times
the second along a fine path in an enlarged fixed collar.
The gradient difference, divided by the input norm, is at most
\begin{equation}\label{cg-schauder-split}
 C\int_0^{r^2}t^{-1+\alpha/2}\,dt
 +Cr\int_{r^2}^1t^{-3/2+\alpha/2}\,dt+Cr
 \le C\left(\frac2\alpha+\frac2{1-\alpha}+1\right)r^\alpha.
\end{equation}
The last term is the massive large-time tail, with a sufficiently
small fixed exponential weight. Larger separations use the
supremum bound. The gradient supremum is integrable for
$\alpha>0$; the field supremum has absolute small-time budget
$t\,t^{-3/2}=t^{-1/2}$. Therefore
\begin{equation}\label{cg-massive-result}
 N_{\alpha,\varpi}(DZ_mD^T\nabla^Tf)
       \le C_\alpha\norm{f}_{\alpha,\varpi}.
\end{equation}
Combining \eqref{cg-remainder} and \eqref{cg-massive-result}
proves the differentiated subtraction estimate. Without the
last input derivative, the absolute budget is
$t\,t^{-(3+\alpha)/2}$, integrable for $\alpha<1$; this gives
$DZD^T:\ell^\infty_\varpi\to N_{\alpha,\varpi}$.
Equation \eqref{cg-constrained} and
\eqref{cg-Gamma-derivative} now prove
Theorem~\ref{cg-schauder}.

\subsection{Full-lattice identities and the application boundary}

The massive estimates use complete fields and global heat-row
cancellation. The remainder estimates use exponentially local
pointwise cell kernels with input derivatives already on the
kernel. No sharp restriction of a H\"older field to one cell is
used. Exact periodization preserves the estimates, including
small compatible tori and their incidence multiplicities.

At each fixed spacing the identities above extend from Hilbert
fields to bounded or admissibly weighted complete fields by
absolute summability of the operator products. In particular
$R=\mathcal LZ\mathcal L$ is exponentially summable; fixed-mesh
equivalence of cell norms suffices for this identity argument.
It is not necessary to call the gauge projection local.
For $\xi\in\mathcal X\cap N_\alpha$, $K\xi$ can be read as
a first-divergence datum with H\"older coefficients, and
\eqref{cg-inverse-identities} agrees with that interpretation.

This supplies the bounded-forcing and equal-exponent jet-forcing
Green inputs for the fixed-flat-gauge nonlinear stationarity
argument. The nonlinear action force, nonlinear constraint
derivatives and strict-minimum argument are separate ingredients.
No moving-gauge nonflat Schauder estimate, mixed retained-history
result or interacting continuum limit is asserted here.
